\section{有限群における同時ブロック置換}
\label{app:group-blocks}

\paragraph{半群型付けの規約．}
群$A$および群への準同型は短縮証明の補助的不変量である．
意味論上の型付けは窓・群積（または局所群型付け）の正語像半群であり，
テンプレート内の空因子の評価には新単位元添加を用いる．
補題~\ref{lem:semigroup-transfer}によって正の長さのタプルの比較，
ならびに今回の特徴正例標本の性質は変化しない．

本付録では定理~\ref{thm:group-typed-shortest}および
補題~\ref{lem:local-group-shortest}に共通する群論的命題を証明する．
可換な場合は零和部分列を用い，非可換な場合はブロックの
相対順序を均質化する．

\begin{definition}[ブロック置換族]
\label{def:block-family}
有限群$A$と$f\ge1$に対し，幅$f$，添字集合$[m]$の
\emph{ブロック置換族}は，$d\le f$本の\emph{列}からなる．
列$c$には左から右へ$\beta_{c,1},\ldots,\beta_{c,q_c}$という
ブロックと群定数$a_{c,0},\ldots,a_{c,q_c}\in A$を並べる．
各ブロック$\beta$は所有者$o(\beta)\in[m]$，
元の値$x_\beta\in A$および置換後の値$x'_\beta\in A$をもち，
各所有者の全列にわたる所有ブロック数は高々$f$である．
$J\subseteq[m]$について
\[
 \mathrm{val}_c(J)=
 a_{c,0}y_{c,1}a_{c,1}\cdots y_{c,q_c}a_{c,q_c},
 \qquad
 y_{c,j}=
 \begin{cases}
 x'_{\beta_{c,j}},&o(\beta_{c,j})\in J,\\
 x_{\beta_{c,j}},&o(\beta_{c,j})\notin J
 \end{cases}
\]
と定める．積の順序は固定し，群の可換性は仮定しない．
\end{definition}

\begin{lemma}[同時ブロック置換]
\label{lem:block-replacement}
幅$f$のブロック置換族について，
(a) $A$が可換で$m>\mathsf d(A^f)$，または
(b) $A$が任意の有限群で$m\ge M_{A,f}$なら，
非空集合$J\subseteq[m]$が存在し，全列$c$に対して
$\mathrm{val}_c(J)=\mathrm{val}_c(\emptyset)$が成立する．
定数$\mathsf d(A^f)$と$M_{A,f}$は
\ref{subsec:group-invariants}節で定義したものを用いる．
\end{lemma}

\subsection{可換群の場合：零和部分列}
\label{app:window-abelian}

\begin{proof}[補題~\ref{lem:block-replacement}(a)の証明]
$A$を加法的に書く．
$\delta_{i,c}$を列$c$のうち所有者$i$に属する全ブロックの
差$x'_\beta-x_\beta$の和とし，
$\delta_i=(\delta_{i,c})_c\in A^d$とおく．
群が可換であるから
$\mathrm{val}_c(J)-\mathrm{val}_c(\emptyset)
 =\sum_{i\in J}\delta_{i,c}$である．
$d\le f$から$\mathsf d(A^d)\le\mathsf d(A^f)$である．
$m>\mathsf d(A^f)$なら列$(\delta_i)_{i=1}^m$は
非空の零和部分列をもつので，該当する所有者集合$J$が
全列の群値を同時に保存する．
\end{proof}



\subsection{一般有限群の場合：順序均質化}
\label{app:window-group}







\begin{proof}[補題~\ref{lem:block-replacement}(b)の証明]
ブロックを一つも所有しない添字があればその添字だけを選べばよい．
以下では各添字が少なくとも一つのブロックを所有するとする．
一つの所有者がもつブロック数・所属列・同一列内の順序で
決まる\emph{単項パターン}は高々
$f\cdot f^f\cdot f!=U_f$種類である．
$m\ge M_{A,f}=U_fN_{A,f}^{2^{f^2}}$なら，
同一パターンの所有者が少なくとも
$N_{A,f}^{2^{f^2}}$個存在する．
それ以外の所有者のブロックは置換せずに定数へ吸収する．

共通パターンのブロックを型$1,\ldots,f'$（$f'\le f$）とし，
それぞれが置かれる列上の位置を点とみなす．
同一列上の型の順序対（同型対を含む）は高々$f^2$個である．
補題~\ref{lem:order-homogenize}を
$N=N_{A,f}=\max\{3,|A|^f\}$として適用すると，
同じパターンをもち，添字間比較がすべて均質な
$N_{A,f}$個の所有者が得られる．
これらについて補題~\ref{lem:group-interval}を適用すれば，
各列の\emph{順序付き}群積を同時に保存する
非空の添字区間が存在する．
\end{proof}

定数$M_{A,f}$は$f$固定では$|A|$の多項式であり，
従来版に現れた多色Ramsey上界より小さい．
ただし$f$への依存は二重指数的に大きい．
